\pdfoutput=1
\documentclass[12pt]{article}

\usepackage{amsmath, amssymb, amsthm}
\usepackage{graphicx}
\usepackage{booktabs}
\usepackage{array}
\usepackage{longtable}
\usepackage{calc}
\usepackage{etoolbox}
\makeatletter
\patchcmd\longtable{\par}{\if@noskipsec\mbox{}\fi\par}{}{}
\makeatother
\IfFileExists{footnotehyper.sty}{\usepackage{footnotehyper}}{\usepackage{footnote}}
\makesavenoteenv{longtable}
\usepackage{xurl}
\usepackage[colorlinks=true, linkcolor=blue, citecolor=blue, urlcolor=blue]{hyperref}
\usepackage{geometry}
\geometry{margin=1in}

\makeatletter
\newsavebox\pandoc@box
\newcommand*\pandocbounded[1]{%
  \sbox\pandoc@box{#1}%
  \Gscale@div\@tempa{\textheight}{\dimexpr\ht\pandoc@box+\dp\pandoc@box\relax}%
  \Gscale@div\@tempb{\linewidth}{\wd\pandoc@box}%
  \ifdim\@tempb\p@<\@tempa\p@\let\@tempa\@tempb\fi
  \ifdim\@tempa\p@<\p@\scalebox{\@tempa}{\usebox\pandoc@box}%
  \else\usebox{\pandoc@box}%
  \fi
}
\makeatother
\providecommand{\tightlist}{%
  \setlength{\itemsep}{0pt}\setlength{\parskip}{0pt}}

\title{Tholonic N-D-C Structure and Its Philosophical Predecessors:\\
\large Spinoza and Leibniz}
\author{J. W. Milton\\[4pt]{\small Clarity Coalition}}
\date{\small 26 June 2026 \\ v1.0}

\begin{document}
\maketitle

\noindent\textbf{Keywords:} tholonic model; Spinoza; Leibniz; substance monism; monadology; N-D-C framework; conatus; pre-established harmony; philosophical predecessors; structural metaphysics; recursive triad

\begin{abstract}
The tholonic N-D-C framework, grounded mathematically in the recursive
iteration of the first three primes and tested empirically across
neural, biological, cosmological, and physical domains in papers
\href{https://github.com/baardev/truevalue/blob/main/docnav/Research/papers/1_recursive-tholonic-five-constants/1_recursive-tholonic-five-constants.pdf}{1}
through
\href{https://github.com/baardev/truevalue/blob/main/docnav/Research/papers/15_consciousness-dmn-d-apparatus/15_consciousness-dmn-d-apparatus.pdf}{15}
of this series, describes how any coherent self-sustaining system
maintains stability through three functional roles: Negotiation (N, the
emergent equilibrium), Definition (D, the constraining parameter), and
Contribution (C, the integrating expressive output). Two
seventeenth-century philosophers, Baruch Spinoza (1632-1677) and
Gottfried Wilhelm Leibniz (1646-1716), independently developed
metaphysical systems that, when examined structurally rather than
doctrinally, anticipate key elements of the tholonic model with
remarkable precision.

This paper performs a systematic structural comparison between the
tholonic model and each philosopher's mature metaphysics. Spinoza's
system contributes three tholonically significant parallels: his
distinction between \emph{Natura naturans} (active, productive Nature)
and \emph{Natura naturata} (the produced modal system) maps onto the
bidirectional N-to-D-and-C recursion; his doctrine of \emph{conatus},
the striving of every mode to persevere in its own being, maps onto the
N-state's structural self-maintenance condition; and the parallelism of
Thought and Extension as attributes of the one substance maps onto D-C
complementarity within a single N level. Leibniz's monadology offers a
tighter structural parallel at the unit level: his monads,
self-contained, self-sustaining, hierarchically nested units that each
mirror the entire universe from their own level, map directly onto
tholons. The doctrine of pre-established harmony, read structurally
rather than theologically, maps onto the emergent balance condition that
produces N from D and C at each hierarchical level.

The paper argues that Spinoza most accurately captured the dynamic
process dimension of the tholonic model, specifically the recursive
productive relation between N and its D-C differentiates, while Leibniz
most accurately captured the structural unit dimension, specifically the
self-similar hierarchically nested tholon. The tholonic model provides a
formal mathematical unification that neither philosopher achieved: a
recursive triadic structure grounded in the first primes that
encompasses both the process insights of Spinoza and the unit insights
of Leibniz. Both philosophers also identified the key failure mode
predicted by the tholonic model, in which a system loses its stable
N-state when D and C fall into severe imbalance, though neither had the
formal vocabulary to name it as such.
\end{abstract}

\section{Introduction}\label{introduction}

\textbf{What this paper provides.} A systematic structural comparison
between the tholonic N-D-C model and the mature metaphysical systems of
Spinoza and Leibniz; identification of which concepts in each
philosopher's system map to which tholonic primitives (N, D, C, tholon,
balance, partial tholon, recursion); an argument for which philosopher
offers the closer structural parallel and why; and an account of what
the tholonic model adds that neither philosopher's system provides.

\textbf{What this paper does not provide.} A comprehensive survey of
either philosopher's complete work; an original contribution to Spinoza
or Leibniz scholarship; a claim that either philosopher consciously
anticipated the tholonic model; or any claim that the historical
priority of these philosophers validates the tholonic model
(philosophical parallel is not scientific evidence). The comparisons are
structural: both systems describe the same underlying pattern from
different starting assumptions.

\textbf{Organization.} §2 briefly recaps the tholonic N-D-C framework
for readers unfamiliar with earlier papers in this series. §3 introduces
the key structural features of Spinoza's system. §4 maps Spinoza's
concepts onto tholonic primitives, identifies the parallels and breaks.
§5 introduces the key structural features of Leibniz's system. §6 maps
Leibniz's concepts onto tholonic primitives. §7 compares the two
philosophers as tholonic precursors and argues for their
complementarity. §8 identifies what the tholonic model adds that neither
philosopher's system can provide. §9 concludes.

\begin{center}\rule{0.5\linewidth}{0.5pt}\end{center}

\section{The Tholonic N-D-C Framework}\label{the-tholonic-n-d-c-framework}

The tholonic model is developed in full in papers
\href{https://github.com/baardev/truevalue/blob/main/docnav/Research/papers/1_recursive-tholonic-five-constants/1_recursive-tholonic-five-constants.pdf}{1}
through
\href{https://github.com/baardev/truevalue/blob/main/docnav/Research/papers/3_minimal-recursive-triadic-framework/3_minimal-recursive-triadic-framework.pdf}{3}
of this series
\href{https://github.com/baardev/truevalue/blob/main/docnav/Research/papers/1_recursive-tholonic-five-constants/1_recursive-tholonic-five-constants.pdf}{Mil26a},
\href{https://github.com/baardev/truevalue/blob/main/docnav/Research/papers/3_minimal-recursive-triadic-framework/3_minimal-recursive-triadic-framework.pdf}{Mil26b},
\href{https://github.com/baardev/truevalue/blob/main/docnav/Research/papers/6_qualitative-nature-integers-triadic-roles/6_qualitative-nature-integers-triadic-roles.pdf}{Mil26c}.
This section provides the minimal recap required for the comparisons
that follow.

\subsection{Core Structure}\label{core-structure}

A \emph{tholon} is any stable, self-sustaining structure that
simultaneously instantiates three functional roles:

\textbf{N (Negotiation)} is the emergent stable equilibrium of the
system: its coherent identity, the state that persists as a consequence
of D and C remaining in approximate balance. N is not any single
component; it is the relational outcome of the D-C dynamic. N does not
exist without both D and C; and D and C do not organize without
producing N.

\textbf{D (Definition)} is the constraining apparatus: all mechanisms
that limit, bound, specify, and filter. D is internally focused. It
determines what the system is, what it excludes, and what boundaries it
maintains. D is the definitional pole: structure, specification,
identity.

\textbf{C (Contribution)} is the integrating-expressive apparatus: all
mechanisms that generate, connect, and output. C is externally focused.
It determines what the system does, what it produces, and how it
connects to the systems around it. C is the contributive pole: flow,
expression, relation.

A \emph{partial tholon} realizes only one or two of these roles and is
structurally unstable. A D-dominant partial tholon collapses into
rigidity: it constrains without producing. A C-dominant partial tholon
dissipates without coherence: it produces without structure. The stable
configuration requires approximate balance, quantified by the B-score:

\[B(D,C) = \frac{2 \cdot \min(D,C)}{D + C} \cdot 100\]

The threshold of structural instability is at \(B = 61.8\)
(\(1/\varphi\)), where \(\varphi\) is the golden ratio.

\subsection{Recursive Self-Similarity}\label{recursive-self-similarity}

The relation between N, D, and C is not a one-way arrow. It is
bidirectional and recursive:

\begin{enumerate}
\def\labelenumi{\arabic{enumi}.}
\tightlist
\item
  A parent N differentiates into D and C, which are the two poles it
  contains.
\item
  D and C negotiate, and their interaction produces a child N.
\item
  The child N becomes the parent N of the next level.
\end{enumerate}

This means N is simultaneously the source from which D and C emerge
\emph{and} the product that D and C jointly produce. The cycle is:

\[\text{Parent } N \;\to\; D + C \;\to\; \text{Child } N \;\to\; \text{Parent } N_{\text{next level}}\]

Five mathematical constants emerge co-emergently from this recursion
when the first three primes (2, 3, 5) are assigned to the N, D, and C
roles: \(\pi/4\), \(\varphi\) (golden ratio), \(e\), \(\sqrt{2}\), and
\(\ln 2\). These constants are not inserted into the model; they arise
from the structure of the recursion itself.

\subsection{Hierarchy}\label{hierarchy}

Every tholon is simultaneously a child N at its own level and a parent N
from which a lower-level D-C differentiation arises. The result is a
hierarchical structure of self-similar units: the smallest tholon embeds
the same triadic pattern as the largest. Each level is structurally
identical to every other level; what varies is the content occupying the
D and C roles.

\subsection{The Partial Tholon and Structural Failure}\label{the-partial-tholon-and-structural-failure}

The tholonic model makes a specific structural prediction about failure:
any system that loses approximate D-C balance will lose its N-state.
This is not a metaphysical claim; it is a structural claim about
stability. A system whose D apparatus dominates becomes rigid and
brittle: it can no longer integrate new C inputs. A system whose C
apparatus dominates becomes incoherent: it can no longer maintain the
boundaries that give it identity. In both cases N dissolves.

This failure mode, and the structural condition for avoiding it, is the
core testable content of the tholonic model.

\begin{center}\rule{0.5\linewidth}{0.5pt}\end{center}

\section{Spinoza's System: Key Structural Features}\label{spinozas-system-key-structural-features}

Baruch Spinoza's mature metaphysics is presented in the \emph{Ethics}
(\emph{Ethica Ordine Geometrico Demonstrata}, 1677), a work organized
using the axiomatic method of Euclidean geometry: definitions, axioms,
propositions, demonstrations, and corollaries {[}Spi77{]}.

\subsection{Substance, Attributes, and Modes}\label{substance-attributes-and-modes}

Spinoza begins with a strict ontological hierarchy:

\textbf{Substance} is that which exists in itself and is conceived
through itself: ``By substance I understand what is in itself and is
conceived through itself: that is, that whose concept does not need the
concept of another thing from which it must be formed'' (E1D3). A
substance is self-sufficient: it requires no external ground for its
existence or intelligibility. Spinoza argues there can be only one such
substance, and it is God, or equivalently, Nature (\emph{Deus sive
Natura}).

\textbf{Attributes} are what the intellect perceives as constituting the
essence of substance (E1D4). Spinoza holds that God/Nature has
infinitely many attributes, but human intellect can access only two:
Thought (\emph{cogitatio}) and Extension (\emph{extensio}). Thought is
the attribute under which the substance expresses itself as minds,
ideas, and consciousness. Extension is the attribute under which it
expresses itself as bodies, matter, and physical reality.

\textbf{Modes} are particular affections of substance: finite, temporal
modifications of the one infinite substance (E1D5). Individual human
bodies, minds, trees, stones, mathematical truths, and emotions are all
modes. A mode requires substance in order to exist and cannot be
conceived without it.

This three-level hierarchy, Substance \textgreater{} Attributes
\textgreater{} Modes, is the backbone of Spinoza's metaphysics. Nothing
exists outside it. Everything that exists is either the one substance,
an attribute of the substance, or a mode of the substance expressed
under one of its attributes.

\subsection{Natura Naturans and Natura Naturata}\label{natura-naturans-and-natura-naturata}

Spinoza makes a key structural distinction within his monism. He
distinguishes between:

\textbf{Natura naturans} (Nature as actively naturing): ``what is in
itself and is conceived through itself, or such attributes of substance
as express eternal and infinite essence, that is, God insofar as He is
considered as a free cause'' (E1P29S). This is the active, productive,
self-causing aspect of the one substance: the infinite power of being
expressing itself.

\textbf{Natura naturata} (Nature as natured, produced): ``everything
that follows from the necessity of God's nature, or from any of God's
attributes, that is, all modes of God's attributes insofar as they are
considered as things that are in God'' (E1P29S). This is the produced
system: the entirety of finite modes, the whole universe of existing
things.

The important structural point is that Natura naturans and Natura
naturata are not two separate entities: they are one substance seen from
two perspectives. Natura naturans is the substance as cause; Natura
naturata is the substance as the totality of its effects. The substance
produces itself, in the sense that everything that exists follows
necessarily from its own nature.

\subsection{Conatus: Structural Self-Preservation}\label{conatus-structural-self-preservation}

In Part III of the \emph{Ethics}, Spinoza introduces one of his most
important and original concepts. Every existing mode, whether physical
or mental, strives to persevere in its own being:

``Each thing, as far as it can by its own power, strives to persevere in
its being'' (E3P6).

The word Spinoza uses is \emph{conatus}, a Latin term meaning effort,
striving, or endeavor. Crucially, Spinoza identifies the conatus of a
mode with its very essence:

``The striving by which each thing strives to persevere in its being is
nothing but the actual essence of the thing'' (E3P7).

The conatus is not an added feature of a mode: it is what the mode
fundamentally is. To be a mode is to be a striving for one's own
continued existence. This applies to every mode without exception:
rocks, plants, animals, ideas, emotions, and human minds all share the
structural property of being constituted by their own self-preserving
activity.

Spinoza argues for conatus through a negative argument: no thing can
destroy itself, because the definition of a thing affirms and does not
deny its own essence (E3P4). What it is to be a cat affirms
cat-existence; it cannot contain a principle of cat-destruction. From
this it follows that each thing must, insofar as it depends on itself
alone, act to preserve its own being.

\subsection{Attribute Parallelism}\label{attribute-parallelism}

Because Thought and Extension are both attributes of the one substance,
everything that occurs in the attribute of Extension has an exact
parallel in the attribute of Thought, and vice versa. A physical event
in a body corresponds precisely to an idea in a mind, not because one
causes the other (Spinoza explicitly rejects inter-attribute causation),
but because they are both expressions of the same modification of the
one substance seen under different attributes. This is the doctrine of
\emph{psychophysical parallelism}.

Spinoza states it directly: ``The order and connection of ideas is the
same as the order and connection of things'' (E2P7). The causal order of
bodies under Extension mirrors the logical order of ideas under Thought,
because both are the same causal order expressed through different
attributes of the same substance.

\begin{center}\rule{0.5\linewidth}{0.5pt}\end{center}

\section{Mapping Spinoza onto the Tholonic Framework}\label{mapping-spinoza-onto-the-tholonic-framework}

\subsection{Substance as N}\label{substance-as-n}

The most immediate structural parallel is between Spinoza's substance
and the tholonic N-state. Both are:

\begin{itemize}
\tightlist
\item
  Self-sustaining (substance ``exists in itself''; N is the emergent
  equilibrium that maintains itself through D-C balance)
\item
  The ground from which everything else derives (modes depend on
  substance; D and C differentiate from N)
\item
  Not any single component but rather the condition for the existence of
  all components (substance is not any particular mode; N is not D or C
  but the equilibrium of both)
\item
  Simultaneously cause and effect within the system (Natura naturans
  produces Natura naturata; parent N produces child N)
\end{itemize}

The parallel is structural, not identical. Spinoza's substance is unique
and infinite: there is exactly one substance, and it is everything. The
tholonic N is level-specific: every tholon has its own N, and N states
are nested at every scale from the smallest tholon to the largest. But
at each level, the N-state has the substance-like property of being the
ground from which that level's D and C emerge and the equilibrium toward
which they tend.

\subsection{Natura Naturans and Natura Naturata as the N-D-C Recursion}\label{natura-naturans-and-natura-naturata-as-the-n-d-c-recursion}

The Natura naturans / Natura naturata distinction is the closest
structural parallel between Spinoza and the tholonic model. Consider the
mapping:

\begin{itemize}
\tightlist
\item
  Natura naturans (the active, self-causing ground of production) = N in
  its role as parent: the differentiated-into-D-and-C source
\item
  Natura naturata (the produced modal system) = D and C in their role as
  the active poles that produce the next level
\item
  The individual mode (the particular thing that strives to persevere) =
  child N: the stable instantiation that emerges from D-C balance at one
  level
\end{itemize}

Spinoza's formulation, that the one substance produces all modes by
necessity from its own nature, mirrors the tholonic cycle: a parent N,
by its own structural necessity, differentiates into D and C, which
negotiate to produce a child N. The substance is both the source and the
totality: it is Natura naturans as the undifferentiated productive
ground, and it is Natura naturata as the totality of everything that has
been produced from it. The tholonic model extends this by making the
recursion explicit at every level and by specifying the D-C balance
condition that the mode must satisfy in order to persist.

Spinoza did not formalize what it means for a mode to ``follow from''
the substance. The tholonic model provides that formalization: a mode is
a stable child N produced by the D-C balance of the level above it.
Modes that cannot achieve D-C balance do not persist: they are the
tholonic equivalent of states that fail to instantiate as stable N.

\begin{figure}
\centering
\pandocbounded{\includegraphics[width=\linewidth,keepaspectratio]{figures/16_philosopher-ndc-triangles.png}}
\caption{Two-panel figure showing the Spinoza-tholon mapping and the
Leibniz-tholon mapping.}
\end{figure}

\subsection{Conatus as the N-State's Self-Maintenance Principle}\label{conatus-as-the-n-states-self-maintenance-principle}

Conatus is, structurally, the most precise tholonic concept in Spinoza.
The tholonic model predicts that every stable N-state tends to maintain
its D-C balance, because it is constituted by that balance: destroy D-C
balance and you destroy N. This is not a value or preference added to
the system; it is a structural consequence of what N is.

This is precisely Spinoza's claim about conatus: the striving to
persevere in being is not an added feature of a mode but its actual
essence. The mode does not have conatus; the mode \emph{is} its conatus.
In tholonic terms: an N-state does not have a tendency to maintain D-C
balance; it \emph{is} the dynamic expression of D-C balance in
equilibrium.

Spinoza also observes that when external causes overcome a mode's
conatus, the mode is destroyed (E3P4's contrapositive: a thing is
destroyed only by external causes). The tholonic parallel: N is
dissolved only when external disruption forces D-C imbalance beyond the
stability threshold. A tholon is not self-destructive; it is destroyed
by being pulled too far from balance by external forces. This is not a
philosophical coincidence; it follows from the same structural logic in
both systems.

\subsection{Attribute Parallelism as D-C Complementarity}\label{attribute-parallelism-as-d-c-complementarity}

Spinoza's two known attributes, Thought and Extension, are not causally
connected (no idea causes a body; no body causes an idea), yet they run
in perfect parallel because they are both expressions of the same
substance. This structure maps onto the tholonic relationship between D
and C within a single N-level:

\begin{itemize}
\tightlist
\item
  Thought (the attribute of minds, ideas, logical structure) maps to D:
  definitional, constraining, formally bounding
\item
  Extension (the attribute of bodies, matter, physical causation) maps
  to C: material, expressive, physically active
\item
  The parallelism (same order in both attributes, same causal structure)
  maps to the tholonic principle that D and C, though functionally
  distinct, are both expressions of the same N-level and must maintain
  proportional balance
\end{itemize}

The tholonic model clarifies why Spinoza's parallelism holds: it holds
because D and C are both constituted by the same N. They are not
independent attributes running in coincidental lockstep; they are the
two poles of a single N-level. Their parallel order is not a mystery; it
is a direct consequence of their common N-ground.

It is important to note, however, that Thought and Extension are not
primary in the tholonic framework. They are specific instantiations of a
more abstract archetype that occupies the D and C roles in any domain.
In the tholonic model, D is most precisely characterized as the role of
\emph{awareness}: the inward-facing, perceiving, defining function that
constitutes the boundary and identity of the system. C is most precisely
characterized as the role of \emph{intention}: the outward-facing,
expressing, and acting function that connects the system to everything
beyond itself. Spinoza's Thought (the capacity to perceive, represent,
and define) is a particular expression of awareness-as-D; his Extension
(the capacity to occupy, resist, and physically act) is a particular
expression of intention-as-C. The tholonic model is domain-neutral
precisely because awareness and intention as archetypes can be
instantiated by entirely different content in different domains: logical
constraint and physical force in Spinoza's ontology, definitional schema
and data throughput in a neural network, regulatory structure and
productive output in a supply chain. The archetype is the same; the
instantiation varies.

\begin{figure}
\centering
\pandocbounded{\includegraphics[width=\linewidth,keepaspectratio]{figures/16_concept-alignment-heatmap.png}}
\caption{Heatmap showing strength of conceptual alignment between each
philosopher's key concepts and the tholonic primitives N, D, C, tholon,
recursion, balance, and partial-tholon.}
\end{figure}

\subsection{Where the Spinoza Mapping Breaks Down}\label{where-the-spinoza-mapping-breaks-down}

The tholonic model and Spinoza's system diverge in three important
respects.

\textbf{Monism vs.~recursive pluralism.} Spinoza's substance is unique:
there is exactly one. The tholonic model is fully hierarchical and
recursive: every tholon is N at its own level and contains D and C that
each produce their own child N. There is no single base tholon; the
recursion extends in both directions. Spinoza's monism forecloses this
recursion by stopping at one substance.

\textbf{Two-attribute limit.} Spinoza holds that human intellect can
access only two attributes: Thought and Extension. The tholonic model
places no type constraint on what occupies the D and C roles, because
those roles are archetypes, not types. The archetypes are awareness (D:
the inward-facing perceiving and defining function) and intention (C:
the outward-facing expressing and acting function). Spinoza's Thought
and Extension are one particular instantiation of these archetypes in
the domain of mind and matter. Any constraining awareness-like function
can occupy D; any expressive intention-like function can occupy C. The
tholonic framework is domain-neutral for this reason: it applies equally
to neural networks, supply chains, atoms, and social institutions,
because awareness and intention as structural archetypes are present in
all of them, even when their specific content looks nothing like
Spinoza's Thought and Extension.

\textbf{Causal necessity and determinism.} Spinoza's system is strictly
determinist: everything follows necessarily from God's nature (E1P29).
The tholonic model does not require determinism. The balance condition
\(D \approx C\) is a structural tendency, not a necessity: systems can
be pushed out of balance and can fail to reach N. The tholonic model
describes attractors, not compelled trajectories.

\begin{center}\rule{0.5\linewidth}{0.5pt}\end{center}

\section{Leibniz's System: Key Structural Features}\label{leibnizs-system-key-structural-features}

Gottfried Wilhelm Leibniz's mature metaphysics is presented most
compactly in the \emph{Monadology} (1714), a work of ninety numbered
propositions that systematically develops the doctrine of monads
{[}Lei14{]}. Supporting detail is found in the \emph{Discourse on
Metaphysics} (1686) {[}Lei86{]}, the \emph{New System} (1695)
{[}Lei95{]}, and the \emph{Theodicy} (1710) {[}Lei10{]}.

\subsection{Monads: Self-Contained Units}\label{monads-self-contained-units}

The fundamental unit of Leibniz's ontology is the monad. Leibniz defines
a monad as a simple substance: it has no parts, no spatial extent, no
internal division (\emph{Monadology} §1). Monads are the true atoms of
nature, but not material atoms: they are metaphysical, not physical.

Each monad is completely self-contained. It has no ``windows'' through
which external influences can enter or internal states can exit
(\emph{Monadology} §7). The states of a monad unfold entirely from
within, according to the monad's own internal law or principle of
activity. Nothing outside a monad can modify it directly; the states it
passes through are entirely determined by its own nature.

Yet each monad expresses, or represents, the entire universe from its
own perspective (\emph{Monadology} §56). Because no monad is entirely
isolated from all others (they are all part of one universe), each monad
contains within it a representation of every other monad. This
representation is not a direct perception; it is a structural
reflection: the monad's internal states correspond to the states of all
other monads at every moment, though the correspondence is more or less
confused and distinct at different levels of the monadic hierarchy.

\subsection{Hierarchy: Bare Monads, Souls, and Rational Monads}\label{hierarchy-bare-monads-souls-and-rational-monads}

Leibniz describes a strict hierarchy of monads organized by the clarity
and distinctness of their perceptions:

\textbf{Bare monads} (entelechies without apperception): the lowest
level. These are the monads that constitute inorganic matter. They have
perceptions, in the sense that they represent the universe, but these
perceptions are entirely confused and indistinct. Bare monads have the
least internal differentiation.

\textbf{Souls} (monads with memory and feeling): the middle level. These
are the monads that constitute animals. They have sentience: they can
remember past states and feel pleasure and pain. Their perceptions are
relatively more distinct than bare monads.

\textbf{Rational souls or spirits} (monads with apperception and
reason): the highest level accessible to creatures. These are human
minds, which can have clear and distinct perceptions and can reason
about necessary truths. They are capable of self-awareness
(\emph{apperception}), which Leibniz distinguishes from mere perception.

\textbf{God} is the supreme monad: infinite, without limitation, whose
perceptions are entirely clear and distinct. God is not a finite monad
among others but the ground of the entire monadic system.

In each living body, one monad is dominant: the \emph{dominant monad}
organizes and integrates the perceptions of the infinitely many
subordinate monads that constitute the body's organs, cells, and
smallest components.

\subsection{Pre-Established Harmony}\label{pre-established-harmony}

Because monads have no windows (they cannot influence each other), the
apparent interaction between them, including the apparent causal
relation between mind and body, requires explanation. Leibniz's
explanation is the doctrine of pre-established harmony
(\emph{Monadology} §78-81).

At creation, God programmed each monad's entire sequence of internal
states such that the internal development of every monad perfectly
corresponds to the internal development of every other monad. The states
of one monad do not cause the states of another monad; they unfold in
parallel, in perfect correspondence, because both sequences were
designed by God to correspond.

For Leibniz this explains the apparent mind-body relationship: my
decision to raise my arm does not cause my arm to rise. Rather, the
mental state corresponding to the decision and the physical state
corresponding to the arm-rising unfold in parallel, in accordance with
their pre-established correspondence. There is no genuine causal
interaction; there is only perfect synchronization.

\subsection{Principles Governing the Monadic System}\label{principles-governing-the-monadic-system}

Leibniz grounds his metaphysics in several structural principles that
govern the behavior of monads:

\textbf{Principle of Sufficient Reason}: nothing is without a sufficient
reason for its being so rather than otherwise (\emph{Monadology} §32).
Every state of every monad has an explanation internal to the monad's
own nature (since it cannot receive external causes).

\textbf{Principle of Continuity} (\emph{lex continuitatis}): nature
makes no jumps. The hierarchy of monads is continuous: between any two
levels of the hierarchy, intermediate levels exist. The gradation from
confused to distinct perception is a continuum, not a step function.

\textbf{Principle of the Identity of Indiscernibles}: no two monads are
exactly alike. Each monad has a unique perspective on the universe, a
unique internal history, and a unique degree of perceptual clarity.

\textbf{Principle of the Best}: God chose to actualize the best of all
possible worlds, where ``best'' is defined by the maximum of perfection
(being, essence, distinctness) given the constraints (\emph{Monadology}
§53-55).

\subsection{Mathematical Grounding}\label{mathematical-grounding}

Leibniz was not only a metaphysician; he was one of the greatest
mathematicians of his era. He independently developed the infinitesimal
calculus (parallel to Newton), invented binary arithmetic, and
contributed foundational work to combinatorics and logic. His
metaphysics and his mathematics were not independent projects; each
informed the other. The monad was in part a metaphysical correlate of
the mathematical point: indivisible, containing information about its
surrounding context, the ultimate atom of analysis. Binary arithmetic,
which Leibniz saw as a representation of creation from nothing (\emph{ex
nihilo}): the generation of all numbers from 0 and 1, directly
influenced his thinking about the monadic hierarchy.

\begin{center}\rule{0.5\linewidth}{0.5pt}\end{center}

\section{Mapping Leibniz onto the Tholonic Framework}\label{mapping-leibniz-onto-the-tholonic-framework}

\subsection{Monads as Tholons}\label{monads-as-tholons}

The monad-tholon correspondence is the closest philosophical parallel to
the tholonic model found in any historical philosopher. The structural
similarities are precise enough to list systematically:

\begin{longtable}[]{@{}
  >{\raggedright\arraybackslash}p{(\linewidth - 2\tabcolsep) * \real{0.5000}}
  >{\raggedright\arraybackslash}p{(\linewidth - 2\tabcolsep) * \real{0.5000}}@{}}
\toprule\noalign{}
\begin{minipage}[b]{\linewidth}\raggedright
Leibniz Monad
\end{minipage} & \begin{minipage}[b]{\linewidth}\raggedright
Tholonic Tholon
\end{minipage} \\
\midrule\noalign{}
\endhead
\bottomrule\noalign{}
\endlastfoot
Simple, self-contained, no internal parts & N-state that constitutes a
level \\
Expresses the entire universe from its perspective & Each tholon embeds
the same triadic pattern as the whole \\
Has an internal principle of activity (entelechy) & N is constituted by
the dynamic D-C balance, not by external imposition \\
Hierarchy from confused to distinct perception & Tholonic hierarchy from
lower to higher levels of N-clarity \\
Dominant monad integrates subordinate monads & Parent N integrates child
N states from the level below \\
No two monads are alike (identity of indiscernibles) & Each tholon has a
unique N-D-C balance configuration \\
\end{longtable}

\begin{figure}
\centering
\pandocbounded{\includegraphics[width=\linewidth,keepaspectratio]{figures/16_monad-tholon-hierarchy.png}}
\caption{Hierarchical nesting diagram comparing the Leibniz monad
hierarchy with the tholonic N-D-C hierarchy.}
\end{figure}

The monad's self-containment and its expression of the whole universe
from its own perspective are both captured by the tholonic principle of
self-similarity: each tholon instantiates the same N-D-C triadic
structure as the whole system, and this self-similarity is why each
tholon can be said to ``mirror'' the structure of the whole. It is not a
poetic metaphor; it is a mathematical consequence of the recursive
triadic structure.

The entelechy, the monad's internal principle of activity that
determines its sequence of states from within, maps onto the emergent
dynamic equilibrium of the N-state. An N-state is not externally driven;
it is constituted by the internal D-C balance. When D and C are in
approximate balance, N exists and is self-sustaining. The tholon's
``activity'' is the continuous dynamic maintenance of D-C balance. This
is not identical to Leibniz's entelechy, but it occupies the same
structural role: the internal source of the unit's activity.

\subsection{Monad Hierarchy as Tholonic Levels}\label{monad-hierarchy-as-tholonic-levels}

Leibniz's three-level hierarchy (bare monad, soul, rational soul) maps
onto the tholonic hierarchy in a direct way. The progression from
confused to distinct perception in Leibniz is a progression in N-clarity
in tholonic terms:

\begin{itemize}
\tightlist
\item
  A bare monad has highly confused, indistinct perceptions: its internal
  D-C balance is minimal or poorly differentiated. It is, in tholonic
  terms, a low-coherence N-state.
\item
  A soul (animal monad) has sentience and memory: a more differentiated
  D (definitional constraints that organize experience over time) and C
  (expressive responses to environment), producing a more coherent N.
\item
  A rational soul has apperception and reason: the most differentiated D
  (logical constraints, necessary truths) and C (rational expression,
  self-awareness), producing the highest achievable N coherence.
\end{itemize}

The dominant monad in a living organism, the one that organizes and
integrates all the subordinate monads, is the tholonic parent N: the
N-state at the highest level of the organism's internal hierarchy. Its
``dominance'' is not a power relation; it is a structural consequence of
being the highest-level integrating tholon of that organism.

\subsection{Pre-Established Harmony as Structural Balance}\label{pre-established-harmony-as-structural-balance}

Leibniz's pre-established harmony is theologically grounded: God
established the correspondence between all monads at creation. The
tholonic model does not require God and does not require an external
establishment of harmony. But the structural function of pre-established
harmony in Leibniz's system, the coordination of all monads without
direct causal interaction, maps onto the tholonic balance condition.

In the tholonic model, D and C do not interact ``directly'' in the sense
of one causing the other through a mechanical transfer. They interact
structurally: the balance between D and C produces N, and N is the
expression of that balance. The ``harmony'' in the tholonic model is not
pre-established by external design; it is the emergent property of a
system that has achieved approximate D-C balance. A tholon that is in
balance is, structurally, a ``harmonious'' monad: its internal states
are coordinated, its expression of the whole is coherent, its N-state is
stable.

The theological scaffolding in Leibniz (God as the establisher of
harmony) is the piece the tholonic model does not replicate. What it
replaces that scaffolding with is the mathematical derivation: the
balance condition \(D \approx C\) is not established by external design
but is the attractor state of any system whose D and C are both active.
Harmony is not given; it is achieved.

\subsection{Windowlessness and Its Tholonic Reinterpretation}\label{windowlessness-and-its-tholonic-reinterpretation}

The most significant structural divergence between Leibniz's monadology
and the tholonic model is the windowlessness doctrine. Leibniz insists
that monads cannot receive external causal influence; all their states
unfold from within. In the tholonic model, tholons explicitly interact
at their D-C boundaries. The D of one tholon can function as C input for
an adjacent tholon, and C outputs from one tholon become N inputs for
the system above it.

Leibniz was driven to windowlessness by his rejection of both Cartesian
interactionism (how can a non-spatial mind move a spatial body?) and
Malebranche's occasionalism (God intervenes to produce every apparent
causal interaction). He had no satisfactory account of how substances
could genuinely transfer causal influence, so he denied that they did.

The tholonic model provides what Leibniz needed but could not construct:
a structural account of inter-level causation. Tholons interact not
through the transfer of ``influx'' (as the scholastic influxus physicus
doctrine had it) but through the structural relation between D-C at one
level and N at the level above. The C output of a child tholon is not
``received'' by the parent tholon through a window; it is structurally
constitutive of the parent's C input. The interaction is hierarchical,
not lateral. Within a level, no monad need transfer influence to
another; across levels, the child N structurally constitutes the
parent's D or C. Leibniz's windowlessness holds within a level; it does
not need to hold across levels.

\subsection{Entelechy, Activity, and the N-State}\label{entelechy-activity-and-the-n-state}

Leibniz's concept of entelechy (from Aristotle: the having of an end
within oneself) is the internal principle of activity that drives each
monad's sequence of states. A monad is not a passive receiver of
impressions; it is an active, self-determining source of its own
development.

In tholonic terms, entelechy corresponds to the N-state's constitutive
dynamic equilibrium. The N-state is not a static point; it is a dynamic
balance. A tholon's ``activity'' is the continuous process of
maintaining D-C equilibrium: adjusting D in response to C pressure, and
C in response to D constraint, in a self-correcting dynamic that keeps N
coherent. This is an internal principle of activity in precisely
Leibniz's sense: it is constituted by the tholon's own D-C structure,
not by anything outside it.

The failure case is also parallel. A monad whose internal development is
somehow disrupted (in Leibniz's terms, by a miracle, an event outside
the natural course) loses the coherence of its perception. A tholon
whose D-C balance is disrupted beyond the stability threshold loses its
N-state. Both systems describe the coherent unit as maintained by
internal activity and dissolved by disruption of that internal
principle.

\subsection{Mathematical Grounding and the Binary Foundation}\label{mathematical-grounding-and-the-binary-foundation}

Leibniz saw binary arithmetic (0 and 1 as the generation of all numbers
from nothing and unity) as a profound metaphysical image of creation. He
also developed the infinitesimal calculus as a tool for describing
continuous change in discrete analytical steps.

The tholonic model's mathematical grounding runs parallel. The
assignment of the first three primes (2, 3, 5) to N, D, and C
respectively, and the derivation of \(\pi\), \(\varphi\), \(e\),
\(\sqrt{2}\), and \(\ln 2\) from the recursive iteration of that
assignment, is the tholonic correlate of Leibniz's binary arithmetic:
both derive fundamental mathematical constants from the simplest
possible combinatorial structure applied recursively. Leibniz chose the
binary pair (0, 1); the tholonic model chooses the triadic prime set (2,
3, 5). Both claim that the mathematical constants are not arbitrary
tools but structural expressions of the simplest possible generating
relations.

\begin{center}\rule{0.5\linewidth}{0.5pt}\end{center}

\section{Comparative Analysis: Which System Is More Structurally Similar?}\label{comparative-analysis-which-system-is-more-structurally-similar}

\subsection{Points of Convergence}\label{points-of-convergence}

Both Spinoza and Leibniz anticipate the tholonic model in distinct and
complementary ways. The table below summarizes the key parallels:

\begin{longtable}[]{@{}
  >{\raggedright\arraybackslash}p{(\linewidth - 4\tabcolsep) * \real{0.3333}}
  >{\raggedright\arraybackslash}p{(\linewidth - 4\tabcolsep) * \real{0.3333}}
  >{\raggedright\arraybackslash}p{(\linewidth - 4\tabcolsep) * \real{0.3333}}@{}}
\toprule\noalign{}
\begin{minipage}[b]{\linewidth}\raggedright
Tholonic Concept
\end{minipage} & \begin{minipage}[b]{\linewidth}\raggedright
Spinoza Parallel
\end{minipage} & \begin{minipage}[b]{\linewidth}\raggedright
Leibniz Parallel
\end{minipage} \\
\midrule\noalign{}
\endhead
\bottomrule\noalign{}
\endlastfoot
N-state (emergent equilibrium) & Substance (self-subsisting ground) &
Monad (self-contained unit) \\
D (Definition/constraint) & Attribute of Thought; limitation within
substance & The monad's confining internal law; structural constraint \\
C (Contribution/expression) & Attribute of Extension; modal expression
outward & The monad's expressive perception; outward mirroring \\
Parent N \(\to\) D+C recursion & Natura naturans differentiating into
modes & Dominant monad differentiating into subordinate monads \\
Child N as stable instantiation & Finite mode persisting through conatus
& Child monad maintaining coherent perception \\
Structural self-preservation & Conatus (the striving of every mode to
persevere) & Entelechy (internal principle of activity) \\
D-C balance condition & Attribute parallelism (Thought = Extension in
structure) & Pre-established harmony (parallel unfolding) \\
Self-similarity across levels & One substance expressed uniformly in all
modes & Every monad mirrors the entire universe \\
Failure mode (N dissolves) & Mode destroyed when conatus overcome by
external causes & Monad destroyed only by miracle (external to its
nature) \\
Mathematical grounding & Geometric axiomatic method; necessary
mathematical truths & Binary arithmetic; calculus; structure of
necessary truths \\
\end{longtable}

\subsection{Which Philosopher Offers the Closer Structural Parallel}\label{which-philosopher-offers-the-closer-structural-parallel}

On balance, Leibniz's system offers the closer structural parallel to
the tholonic model. Three considerations support this:

\textbf{Unit correspondence.} The monad-tholon correspondence is precise
at the level of the individual unit. Both are self-contained,
self-sustaining, hierarchically organized, and self-similar. Spinoza's
modes are structurally comparable but are entirely dependent on the one
substance; they do not have the self-contained, self-mirroring quality
of monads. The tholonic model is fundamentally about units (tholons),
not about a single infinite ground.

\textbf{Hierarchical structure.} Leibniz's three-level hierarchy (bare
monad, soul, rational soul) with dominant and subordinate monads maps
precisely onto the tholonic hierarchy of parent N, child N, and
subordinate tholons. Spinoza's system has a hierarchy (substance,
attributes, modes), but the attributes are not organized hierarchically
in the same way; they run in parallel rather than nesting recursively.

\textbf{Mathematical ambition.} Leibniz's binary arithmetic and calculus
reflect the same ambition as the tholonic model: to derive fundamental
structural properties from the simplest mathematical operations applied
recursively. Spinoza's geometric method is structurally analogous in
form (deduction from axioms) but does not make the same claim about
mathematical constants emerging from the system's recursive structure.

\subsection{Where Spinoza Is Irreplaceable}\label{where-spinoza-is-irreplaceable}

Despite the closer structural parallel with Leibniz, Spinoza captures
something that Leibniz does not: the \emph{dynamic process} dimension of
the N-D-C recursion.

The Natura naturans / Natura naturata distinction is the most direct
structural parallel to the bidirectional recursion in the tholonic model
(parent N produces D and C; D and C produce child N). Leibniz's
pre-established harmony, by contrast, is static: the coordination was
established once, at creation, and does not require ongoing dynamic
production. In the tholonic model, N is dynamically constituted at every
moment by the active balance of D and C. Spinoza's substance is
productive (it actively produces its modes) rather than pre-harmonized.
This dynamic quality of Spinoza's system is a closer match to the
tholonic recursion.

Similarly, conatus is a richer account of the N-state's self-maintenance
than Leibniz's entelechy. Conatus is not merely an internal principle;
it is explicitly a \emph{striving}, an active resistance to dissolution.
This maps onto the tholonic principle that a stable N actively maintains
D-C balance, adjusting when disturbed. Leibniz's entelechy is more
passive: it determines the monad's sequence of states without specifying
that the monad actively resists external disruption.

\subsection{Complementarity}\label{complementarity}

Spinoza and Leibniz are best understood, from the tholonic perspective,
as complementary partial captures of the tholonic structure:

\begin{itemize}
\tightlist
\item
  Spinoza captures the \textbf{process} dimension: the recursive
  productive relation between N and its D-C differentiates, and the
  active self-maintenance of each instantiated state (conatus).
\item
  Leibniz captures the \textbf{unit} dimension: the self-contained
  self-sustaining tholon, the hierarchical nesting, and the structural
  correspondence between levels (pre-established harmony as structural
  balance).
\end{itemize}

Neither philosopher reached the tholonic synthesis because neither had
the formal vocabulary to express the balance condition (\(D \approx C\)
as the stability criterion) or the recursive structure (N simultaneously
as source and product). These are the formal additions that the tholonic
model supplies.

\begin{center}\rule{0.5\linewidth}{0.5pt}\end{center}

\section{What the Tholonic Model Adds}\label{what-the-tholonic-model-adds}

The comparison clarifies not only what Spinoza and Leibniz anticipated
but also what they were unable to provide. The tholonic model adds four
structural elements that neither philosopher's system contains:

\textbf{1. The balance condition.} Neither Spinoza nor Leibniz specifies
a quantitative balance condition between the constraining and expressive
poles of their systems. The tholonic model specifies this precisely:
\(B = 2 \cdot \min(D,C)/(D+C) \cdot 100\), with a stability threshold at
\(61.8\) (\(1/\varphi\)). This transforms philosophical intuition into a
testable, measurable criterion.

\textbf{2. The partial tholon and its failure modes.} Both philosophers
implicitly recognize failure modes: for Spinoza, a mode is destroyed
when external causes overcome its conatus; for Leibniz, a monad's
coherence is disrupted by miracle. But neither has a formal account of
\emph{how} this failure occurs structurally. The tholonic model
specifies two failure modes (D-dominance: rigidity, brittleness;
C-dominance: dissipation, incoherence) and predicts when each will occur
based on the D-C ratio.

\textbf{3. The derivation of mathematical constants.} Leibniz recognized
a connection between mathematical structure and metaphysical structure
(binary arithmetic and the monadic system), but did not derive the
mathematical constants of his system from the triadic recursion. The
tholonic model derives \(\pi\), \(\varphi\), \(e\), \(\sqrt{2}\), and
\(\ln 2\) directly from the recursive iteration of the first three
primes assigned to N, D, and C. These constants are not philosophical
symbols; they are load-bearing structural elements that appear
empirically in the domains where the tholonic model has been tested
(neural network architecture, atomic structure, cosmological phase
transitions, cancer progression).

\textbf{4. Domain-neutral applicability.} Spinoza's system is
constrained to the two attributes knowable to human intellect (Thought
and Extension). Leibniz's system is constrained by its theological
framework (God as supreme monad, pre-established harmony as divine
design). The tholonic model is domain-neutral: D and C can be occupied
by any constraining and expressive function in any domain. This
generality is what allows the same formal structure to describe supply
chains, neural networks, biological organisms, atoms, and cosmological
phase transitions within a single framework.

\begin{center}\rule{0.5\linewidth}{0.5pt}\end{center}

\section{Conclusions}\label{conclusions}

Two of the most rigorous metaphysical thinkers of the early modern
period independently developed systems that anticipate the tholonic
N-D-C framework. Spinoza's \emph{Ethics} captures the recursive
productive process at the core of the tholonic model: the active
differentiation of N into D and C (Natura naturans), the striving of
each produced state to maintain itself (conatus), and the parallelism of
the constraining and expressive poles within a shared ground (attribute
parallelism). Leibniz's \emph{Monadology} captures the structural unit
at the core of the tholonic model: the self-contained hierarchically
nested tholon (monad), the self-similarity across levels (each monad
mirrors the whole), and the structural balance that allows each unit to
maintain its coherence (pre-established harmony as emergent
coordination).

Leibniz's system offers the closer overall structural parallel to the
tholonic model, because the tholonic model is fundamentally a theory of
structural units rather than of a single infinite substrate, and because
Leibniz's hierarchical monad structure maps precisely onto the tholonic
recursive hierarchy. Spinoza's system is the closer parallel
specifically for the process dimension: the recursive productive
relation between N and D+C, and the active character of structural
self-maintenance.

The tholonic model adds what neither philosopher could provide: the
formal balance condition, the two failure-mode predictions, the
derivation of the five mathematical constants, and the domain-neutral
applicability that allows the same structure to be tested and measured
across radically different physical and informational domains. Both
Spinoza and Leibniz were, in structural terms, doing tholonic analysis
with the conceptual vocabulary available in the seventeenth century. The
tholonic model provides the formalization they were reaching toward.

\begin{center}\rule{0.5\linewidth}{0.5pt}\end{center}

\section{References}\label{references}

{[}Spi77{]} Spinoza, B. \emph{Ethics (Ethica Ordine Geometrico
Demonstrata).} 1677. Trans. E. M. Curley. \emph{The Collected Works of
Spinoza.} Princeton University Press, 1985.

{[}Lei14{]} Leibniz, G. W. \emph{The Monadology.} 1714. Trans. R. Ariew
and D. Garber. \emph{Philosophical Essays.} Hackett, 1989.

{[}Lei86{]} Leibniz, G. W. \emph{Discourse on Metaphysics.} 1686. Trans.
R. Ariew and D. Garber. \emph{Philosophical Essays.} Hackett, 1989.

{[}Lei95{]} Leibniz, G. W. \emph{New System of the Nature and
Communication of Substances.} 1695. Trans. R. Ariew and D. Garber.
\emph{Philosophical Essays.} Hackett, 1989.

{[}Lei10{]} Leibniz, G. W. \emph{Theodicy.} 1710. Trans. E. M. Huggard.
Open Court, 1952.

\href{https://github.com/baardev/truevalue/blob/main/docnav/Research/papers/1_recursive-tholonic-five-constants/1_recursive-tholonic-five-constants.pdf}{Mil26a}
Milton, J. W. \emph{Recursive Tholonic Five Constants.} Clarity
Coalition,
\href{https://github.com/baardev/truevalue/blob/main/docnav/Research/papers/1_recursive-tholonic-five-constants/1_recursive-tholonic-five-constants.pdf}{paper
1}, 2026.

\href{https://github.com/baardev/truevalue/blob/main/docnav/Research/papers/3_minimal-recursive-triadic-framework/3_minimal-recursive-triadic-framework.pdf}{Mil26b}
Milton, J. W. \emph{Minimal Recursive Triadic Framework.} Clarity
Coalition,
\href{https://github.com/baardev/truevalue/blob/main/docnav/Research/papers/3_minimal-recursive-triadic-framework/3_minimal-recursive-triadic-framework.pdf}{paper
3}, 2026.

\href{https://github.com/baardev/truevalue/blob/main/docnav/Research/papers/6_qualitative-nature-integers-triadic-roles/6_qualitative-nature-integers-triadic-roles.pdf}{Mil26c}
Milton, J. W. \emph{Qualitative Nature of Integers and Triadic Roles.}
Clarity Coalition,
\href{https://github.com/baardev/truevalue/blob/main/docnav/Research/papers/6_qualitative-nature-integers-triadic-roles/6_qualitative-nature-integers-triadic-roles.pdf}{paper
6}, 2026.

\href{https://github.com/baardev/truevalue/blob/main/docnav/Research/papers/9_tholonic-model-standard-model-alternative/9_tholonic-model-standard-model-alternative.pdf}{Mil26d}
Milton, J. W. \emph{Tholonic Model and the Standard Model: An
Alternative Framework.} Clarity Coalition,
\href{https://github.com/baardev/truevalue/blob/main/docnav/Research/papers/9_tholonic-model-standard-model-alternative/9_tholonic-model-standard-model-alternative.pdf}{paper
9}, 2026.

\href{https://github.com/baardev/truevalue/blob/main/docnav/Research/papers/10_tholonic-neural-architecture/10_tholonic-neural-architecture.pdf}{Mil26e}
Milton, J. W. \emph{Neural Networks as Tholonic Systems.} Clarity
Coalition,
\href{https://github.com/baardev/truevalue/blob/main/docnav/Research/papers/10_tholonic-neural-architecture/10_tholonic-neural-architecture.pdf}{paper
10}, 2026.

\href{https://github.com/baardev/truevalue/blob/main/docnav/Research/papers/12_newton-tholonic-framework/12_newton-tholonic-framework.pdf}{Mil26f}
Milton, J. W. \emph{Newton and the Tholonic Framework.} Clarity
Coalition,
\href{https://github.com/baardev/truevalue/blob/main/docnav/Research/papers/12_newton-tholonic-framework/12_newton-tholonic-framework.pdf}{paper
12}, 2026.

\href{https://github.com/baardev/truevalue/blob/main/docnav/Research/papers/13_dual-entropy-tholonic-cosmology/13_dual-entropy-tholonic-cosmology.pdf}{Mil26g}
Milton, J. W. \emph{Dual Entropy Tholonic Cosmology.} Clarity Coalition,
\href{https://github.com/baardev/truevalue/blob/main/docnav/Research/papers/13_dual-entropy-tholonic-cosmology/13_dual-entropy-tholonic-cosmology.pdf}{paper
13}, 2026.

\href{https://github.com/baardev/truevalue/blob/main/docnav/Research/papers/14_cancer-d-collapse/14_cancer-d-collapse.pdf}{Mil26h}
Milton, J. W. \emph{Cancer and D-Collapse.} Clarity Coalition,
\href{https://github.com/baardev/truevalue/blob/main/docnav/Research/papers/14_cancer-d-collapse/14_cancer-d-collapse.pdf}{paper
14}, 2026.

\href{https://github.com/baardev/truevalue/blob/main/docnav/Research/papers/15_consciousness-dmn-d-apparatus/15_consciousness-dmn-d-apparatus.pdf}{Mil26i}
Milton, J. W. \emph{The Default Mode Network as D-Apparatus.} Clarity
Coalition,
\href{https://github.com/baardev/truevalue/blob/main/docnav/Research/papers/15_consciousness-dmn-d-apparatus/15_consciousness-dmn-d-apparatus.pdf}{paper
15}, 2026.

{[}Gar02{]} Garrett, D. ``Spinoza's Conatus Argument.'' In
\emph{Spinoza: Metaphysical Themes.} ed.~Koistinen and Biro. Oxford
University Press, 2002.

{[}DelR96{]} Della Rocca, M. \emph{Representation and the Mind-Body
Problem in Spinoza.} Oxford University Press, 1996.

{[}Rut95{]} Rutherford, D. \emph{Leibniz and the Rational Order of
Nature.} Cambridge University Press, 1995.

{[}Loe69{]} Loeb, L. \emph{From Descartes to Hume: Continental
Metaphysics and the Development of Modern Philosophy.} Cornell
University Press, 1981.

{[}Ish06{]} Ishiguro, H. \emph{Leibniz's Philosophy of Logic and
Language.} Cambridge University Press, 1990.

\end{document}
